\documentclass[11pt]{article}
\usepackage[margin=1in]{geometry}
\usepackage[T1]{fontenc}
\usepackage{lmodern,amsmath,amssymb,amsthm,mathtools,hyperref}
\setlength{\emergencystretch}{3em}
\newtheorem{theorem}{Theorem}[section]
\newtheorem{lemma}[theorem]{Lemma}
\newtheorem{corollary}[theorem]{Corollary}
\newtheorem{proposition}[theorem]{Proposition}
\newcommand{\E}{\mathbb E}
\newcommand{\tr}{\operatorname{tr}}
\title{Four-phase constructions, weighing completion,\\and exact Gram-defect repair}
\author{Joshua Mark\thanks{ORCID: \href{https://orcid.org/0009-0003-5905-2382}{0009-0003-5905-2382}. AI tools assisted with derivation, independent review, computation and exposition.}}
\date{20 November 2026 --- research draft; priority not established}
\begin{document}\maketitle

\begin{abstract}
Four formulations of the real Hadamard existence problem are examined.
For complex Hadamard matrices, the operator norm of the entrywise fourth
power exactly detects conversion to the quaternary phase alphabet and
gives a quantitative obstruction for Fourier-block families.
For weighing matrices with two zeros per row, all zero-preserving
Hadamard completions reduce to eight root choices; an explicit
order-twelve example has no real completion supported on its zeros.
Rank-two Gram errors have an integer-quantized energy spectrum, and
the energy $8(n-1)$ alone certifies a unique one-entry Hadamard repair.
Finally, flat quantum measurements give exact commuting criteria,
but even flat dilations with vanishing relative dimension overhead
do not ensure nearby projective descent with fixed sign labels.
These results provide local certificates and obstructions, not an
all-order construction. Bibliographic priority is not established.
\end{abstract}
The full conjecture asks for a real sign matrix $H$ with
$HH^T=nI_n$ for every positive $n$ divisible by four.
Each section below keeps its additional hypotheses explicit.
The finite checks certify examples; the general statements have paper proofs.
\section{A four-phase route and a quantitative Fourier obstruction}
\label{sec:four-phase-route}

A quaternary Hadamard matrix $B$ of order $m$, also denoted
$\mathrm{BH}(m,4)$, has entries in $\mu_4=\{1,-1,i,-i\}$ and
$BB^*=mI_m$. Such matrices give real Hadamard matrices of order
$2m$. This is a classical construction, not a new existence theorem;
see Butson, \emph{Generalized Hadamard matrices} (1962),
\url{https://doi.org/10.1090/S0002-9939-1962-0142557-0},
and Egan et al., the morphism discussion in
\url{https://ajc.maths.uq.edu.au/pdf/73/ajc_v73_p501.pdf}.

This route permits arbitrary complex phases as an initial search
space, which is nonempty in every dimension, but the final four-phase
constraint must actually be met. The results below give an exact
phase criterion and a quantitative obstruction to some obvious
Fourier deformations. Priority for the combined observations remains
unestablished; neither gives an all-order existence theorem.

\subsection{Exact realification and the missing monomial structure}
For a complex matrix $Z$ write $\mathcal R(Z)$ for its realification,
replacing $x+iy$ by $\left(\begin{smallmatrix}x&-y\\y&x\end{smallmatrix}\right)$.
It respects products and sends conjugate transpose to real transpose.
For $B\in\mathrm{BH}(m,4)$, put
\[
                   H=\mathcal R((1+i)B).
\]
Its entries are signs and $HH^T=2mI_{2m}$. Thus proving
$\mathrm{BH}(2t,4)$ existence for every positive $t$ would suffice
for the full real Hadamard conjecture. No converse or proof of this
stronger target is being assumed.

\begin{proposition}[A signed complex-structure criterion]
\label{quat:structure}
A real Hadamard matrix $H$ of order $2m$ is equivalent, by row/column
permutations and sign changes, to such a realification if and only if there are signed
permutation matrices $J_{\rm row},J_{\rm col}$ with
\[
 J_{\rm row}^2=J_{\rm col}^2=-I,\qquad
                         HJ_{\rm col}=J_{\rm row}H.
\]
\end{proposition}
\begin{proof}
Each signed permutation squaring to $-I$ pairs coordinates, with
opposite signs on each two-cycle. Signed permutations of coordinates
put each into the canonical block form
$J=\operatorname{diag}(\left(\begin{smallmatrix}0&-1\\1&0\end{smallmatrix}\right))$.
Transform rows and columns independently. The intertwining equation
becomes $H'J=JH'$, so every two-square block is
$\left(\begin{smallmatrix}a&-b\\b&a\end{smallmatrix}\right)$ with
$a,b\in\{\pm1\}$. Hence $H'=\mathcal R(Z)$ and
$ZZ^*=2mI_m$. Dividing $Z$ by $1+i$ gives a quaternary matrix
of Gram $mI_m$. The converse follows by taking the canonical $J$
on both sides of the displayed realification.
\end{proof}

The signed-permutation condition is essential. For any even-order
scaled orthogonal $H$ and any orthogonal complex structure $J_{\rm col}$,
one can set $J_{\rm row}=HJ_{\rm col}H^T/(2m)$. This always gives
an intertwiner and a real orthogonal complex structure, but usually
not a signed permutation. It need not pair individual coordinates,
so it does not give a four-phase matrix. The checker verifies this
failure for the default adjacent pairing on a normalized Paley $H_{12}$;
it does not claim that no other pairing works.

\subsection{An invariant that detects four-phase accessibility}
Let $F$ be any complex Hadamard matrix of order $n$ with unit-modulus
entries. Write $F^{\circ4}$ for its entrywise fourth power and define
\[
                 \rho_4(F)=\frac{\|F^{\circ4}\|_{\rm op}}n.
\]
Multiplying rows or columns by unit phases, or permuting them, does
not change $\rho_4$, since these operations induce unitary monomial
operations on $F^{\circ4}$.

\begin{theorem}[Exact criterion and quantitative separation]
\label{quat:phase-defect}
One has $\rho_4(F)\le1$, with equality if and only if $F$ can be
converted into a quaternary Hadamard matrix by row and column phases.
For every matrix $F'$ equivalent to $F$ under row/column phases and
permutations, and every matrix $Q$ with entries in $\mu_4$,
\[
              \frac{\|F'-Q\|_F^2}{n^2}
                    \ge\frac{1-\rho_4(F)}8.
\]
The comparison matrix $Q$ need not be Hadamard.
\end{theorem}
\begin{proof}
The Frobenius norm of $M=F^{\circ4}$ is $n$, so its operator norm
is at most $n$. Equality is equivalent to $M$ having rank one.
A rank-one matrix with unit-modulus entries factors as $M_{ij}=u_i v_j$
with unit $u_i,v_j$; choose fourth-root row and column phases to
remove them. The transformed entries then have fourth power one,
and the Hadamard Gram is preserved. The reverse implication is immediate.

For the quantitative bound, set $M'=(F')^{\circ4}$ and let $J_n$
be the all-one matrix. Since $|z^4-w^4|\le4|z-w|$ on the unit circle,
\[
 16\|F'-Q\|_F^2\ge\|M'-J_n\|_F^2
 =2n^2-2\operatorname{Re}(\mathbf1^TM'\mathbf1)
 \ge2n^2-2n\|M'\|_{\rm op}.
\]
Use invariance of the operator norm under the phase and permutation
changes to obtain the result.
\end{proof}

This criterion is a spectral version of the classical four-cycle
phase invariants of complex Hadamard equivalence. The relevant
background and the block deformation construction below are discussed
in Tadej--\.Zyczkowski, \emph{A concise guide to complex Hadamard matrices},
\url{https://arxiv.org/abs/quant-ph/0512154}.
It does not show that an optimizer of $\rho_4$ reaches one.

\subsection{Fourier blocks retain a definite obstruction}
Let $\mathcal F_m=(\zeta_m^{ij})_{0\le i,j<m}$ and $d=\gcd(m,4)$.
Direct character orthogonality gives
\[
 (\mathcal F_m^{\circ4}(\mathcal F_m^{\circ4})^*)_{ij}
 =\begin{cases}m&m\mid4(i-j),\\0&\text{otherwise}.\end{cases}
\]
There are $m/d$ all-one blocks of size $d$ after row reordering, so
\[
 \operatorname{rank}(\mathcal F_m^{\circ4})=m/d,\qquad
 \rho_4(\mathcal F_m)=\sqrt{d/m}.
\]
In particular Fourier row/column phases can produce a quaternary
matrix exactly when $m$ divides four. Arbitrary phases do not remove
this obstruction.

More generally take an order-$r$ complex Hadamard matrix $A$,
unit phases $u_{bi}$, and the classical block family
\[
 F_{(a,i),(b,j)}=A_{ab}u_{bi}\zeta_m^{ij}.
\]
It is Hadamard: inner products vanish in the Fourier factor for
distinct $i$, and otherwise reduce to those of $A$.
Every fourth-power block has operator norm $\sqrt{md}$.
For a vector partitioned into $r$ blocks, the triangle and
Cauchy--Schwarz inequalities therefore bound the full operator norm
by $r\sqrt{md}$. Consequently
\[
 \rho_4(F)\le\sqrt{d/m},\qquad
 \inf_{F'\sim F,\ Q\in\mu_4^{rm\times rm}}
 \frac{\|F'-Q\|_F^2}{(rm)^2}
 \ge\frac{1-\sqrt{d/m}}8.
\]
Thus these Fourier-block deformations cannot reach the four-phase
target when $m\nmid4$, even with arbitrarily chosen deformation phases.
This excludes one search family, not every continuous path through
complex Hadamard matrices.

At order six the distinction is concrete: Fourier $\mathcal F_6$
has $\rho_4=1/\sqrt3<1$, whereas a symmetric conference matrix $C$
of order six gives $B=I+iC\in\mathrm{BH}(6,4)$, with $\rho_4(B)=1$.
Its realification gives the classical $H_{12}$. The checker verifies
this and the corresponding orders fourteen and eighteen in exact
Gaussian-integer arithmetic, producing real orders twenty-eight
and thirty-six. The conference construction comes from Paley's
known prime-field matrices, not an unknown-order existence argument.

The sufficient research target is now precise: find a four-phase
complex Hadamard matrix at every even order, or force $\rho_4=1$
somewhere in each corresponding complex-Hadamard space.
Those spaces are nonempty, but their Fourier phase or block
families alone cannot supply the general conclusion.
\section{Two-zero weighing matrices: an eight-choice algorithm and a real obstruction}
\label{sec:weighing-fill}

A weighing matrix $W(n,k)$ has entries $0,\pm1$ and satisfies
$WW^{\mathsf T}=kI_n$. We ask whether its zero entries can be
filled by signs, preserving every nonzero entry, to give a Hadamard
matrix of the \emph{same} order. This constraint is stronger than
merely knowing that the order admits some Hadamard matrix.
The orthogonal-design viewpoint is classical; see Geramita--Verner,
\emph{Orthogonal Designs with Zero Diagonal}, Theorems 1--2,
\url{https://doi.org/10.4153/CJM-1976-028-3}.
The structural proofs and explicit finite certificate below are
independently derived. No bibliographic priority is claimed.

\subsection{One zero: a signed-cycle criterion and the classical positive case}
In $W(n,n-1)$ each row and column has exactly one zero, so a column
permutation puts all zeros on the diagonal. A filling has the form
$W+\operatorname{diag}(d_i)$, with $d_i\in\{\pm1\}$. Its Gram
matrix is $nI$ exactly when
\[
 d_id_j=-w_{ij}w_{ji}\quad(i\ne j).
\]
Thus completion is equivalent to the complete signed graph with edge
labels $s_{ij}=-w_{ij}w_{ji}$ being balanced: the product around every
triangle is $1$. Taking $d_0=1$ determines all other signs
$d_i=s_{0i}$; check the remaining edges. There are either zero or
two fillings, the latter differing by global negation of the fill.
This gives an $O(n^2)$ test and a triangle certificate on failure.

For conference matrices this reproduces a classical theorem:
at $n\equiv0\pmod4$, row and column sign changes make the matrix
skew-symmetric, so diagonal filling always succeeds. At
$n\equiv2\pmod4$, $n>2$, the matrix is similarly equivalent to a
symmetric conference matrix and same-order Hadamard filling is
impossible. These normalization statements are the
Delsarte--Goethals--Seidel result restated as Theorem 2 in the cited
primary article, not a new existence assertion.

\subsection{Two zeros force much more structure than arbitrary cycles}
Put
\[
 F=\begin{pmatrix}1&1\\1&-1\end{pmatrix},\qquad
 \mathcal D=\{2\times2\text{ signed permutation matrices}\},
 \qquad|\mathcal D|=8.
\]

\begin{theorem}[Forced block structure]
Let $4\mid n$ and $W$ be $W(n,n-2)$. After independently permuting
rows and columns, its zeros are precisely the diagonal $2\times2$
blocks. Every off-diagonal $2\times2$ block is itself Hadamard.
Equivalently, with $m=n/2$,
\[
 A=\tfrac12W(I_m\otimes F)
\]
has zero diagonal blocks and off-diagonal blocks $A_{ij}\in\mathcal D$,
and $AA^{\mathsf T}=(m-1)I_n$.
\end{theorem}
\begin{proof}
Two row zero-neighborhoods each have size two. If their intersection
has size $t$, their nonzero overlap has size $n-4+t$, which must be
even for their sign products to sum to zero. Hence $t$ is zero or
two. Each column has exactly two zeros, so rows with the same
two-element neighborhood occur in pairs, and distinct neighborhoods
are disjoint. The zero support is a union of $K_{2,2}$ components,
not longer bipartite cycles. This proves the block normalization.

Every row of an off-diagonal block $W_{ij}$ consists of two signs.
Multiplication by $F/2$ turns it into a row with one nonzero sign,
so $A$ is integral, with one nonzero entry per row of each
off-diagonal block. Also $AA^{\mathsf T}=(m-1)I_n$.
As $m$ is even, reduction modulo two makes $A$ orthogonal and
invertible. Let $L$ be the subspace of vectors constant on each
coordinate pair. The action of $A$ on $L$ is $J_m-I_m$, whose
square is $I_m$ over $\mathbb F_2$ because $m$ is even.
Consequently $A^{\mathsf T}|_L=A^{-1}|_L=J_m-I_m$ as well.
Applied to a pair indicator, this says that the two column sums
of every off-diagonal block are both odd. Since that block already
has exactly one nonzero entry in each row, it must have exactly
one in each column too. It is therefore a signed permutation.
Finally $W_{ij}=A_{ij}F$ is Hadamard.
\end{proof}

\begin{theorem}[All sign fillings in quadratic time]
For a normalized $W(n,n-2)$, all zero-preserving Hadamard fillings
can be listed by testing eight initial states. There are at most
eight fillings, and the test uses $O(n^2)$ integer operations.
Specifically, a filling has diagonal blocks $S_i=R_iF$,
$R_i\in\mathcal D$. Choose $R_0\in\mathcal D$, force
\[
 R_j=-A_{j0}R_0^{\mathsf T}A_{0j}\quad(j>0),
\]
and accept exactly when
\[
 A_{ij}R_j^{\mathsf T}+R_iA_{ji}^{\mathsf T}=0
 \quad\text{for every }i<j.
\]
Each accepted state gives one complete Hadamard matrix.
\end{theorem}
\begin{proof}
Let $X$ be supported on the zero blocks. Inside each row pair,
the cross terms $WX^{\mathsf T}+XW^{\mathsf T}$ vanish, so the
Hadamard equations force $S_iS_i^{\mathsf T}=2I_2$. Thus $S_i$
is a two-square Hadamard matrix and $R_i=S_iF/2$ is a signed
permutation. Different pairs have $XX^{\mathsf T}=0$, leaving
exactly the displayed cross equations. The equations with pair
$0,j$ solve uniquely for $R_j$ as shown, since $A_{0j}$ and
$A_{j0}$ are orthogonal. All forced states remain in $\mathcal D$.
Testing the other pairs is necessary and sufficient, and all
fillings arise this way. There are eight roots and $O(n^2)$
constant-size block operations per root. The order-four example
below has eight fillings, so the cardinality bound is sharp.
\end{proof}

Thus the general two-zero problem is not an unconstrained
exponential search over the $2n$ missing signs. Its obstruction
is a small matrix consistency condition around block triangles.
Nevertheless consistency can fail, even over the reals.

\subsection{An explicit $W(12,10)$ that cannot be filled even by real entries}
Let $I$ be the two-square identity and put
\[
 D=\begin{pmatrix}1&0\\0&-1\end{pmatrix},\quad
 S=\begin{pmatrix}0&1\\1&0\end{pmatrix},\quad
 J=DS=\begin{pmatrix}0&1\\-1&0\end{pmatrix}.
\]
Define the following six-square \emph{block} matrix, whose entries
are two-square matrices:
\[
 A=\begin{pmatrix}
0&I&D&D&D&D\\
S&0&-I&-I&I&I\\
S&-D&0&I&-D&D\\
S&-D&I&0&D&-D\\
S&D&-D&D&0&-I\\
S&D&D&-D&-I&0
\end{pmatrix},\qquad W=A(I_6\otimes F).
\]
Direct integer block multiplication gives $AA^{\mathsf T}=5I_{12}$,
so $WW^{\mathsf T}=10I_{12}$. Its nonzero entries are signs and
its zeros are exactly the diagonal two-square blocks.

Suppose first that even \emph{real} zero replacements could give
a matrix $H$ with $HH^{\mathsf T}=12I_{12}$. As in the preceding
proof the fill blocks have $S_iS_i^{\mathsf T}=2I$, so
$R_i=S_iF/2\in O(2)$, and the same propagation formula applies.
Writing $R=R_0$, the specified blocks give
\[
 R_1=-SR^{\mathsf T},\qquad
 R_2=R_4=-SR^{\mathsf T}D.
\]
The pair $1,2$ cross equation becomes
$DRS+SR^{\mathsf T}D=0$, equivalently
\[
 R+JR^{\mathsf T}J=0.
\]
For every real two-square matrix,
$JR^{\mathsf T}J=R-\operatorname{tr}(R)I$, so $R$ must be
scalar. Its orthogonality then forces $R=\lambda I$ with
$\lambda=\pm1$. But pair $2,4$ requires
$RS+SR^{\mathsf T}=0$, or $2\lambda S=0$, a contradiction.
This proves that no real orthogonal completion can change only
the zero entries, and in particular no sign filling exists.
Hadamard matrices of order twelve do exist; the obstruction
concerns this prescribed set of nonzero entries.

The checker embeds the full integer $W$, verifies its Gram matrix,
and records a failing block pair for each of the eight root states.
Six states already fail on triangle $0,1,2$; the remaining scalar
states fail on triangle $0,2,4$. The symbolic argument above
certifies the stronger real obstruction without numerical optimization.

\subsection{Positive cases and the limit of the route}
For an order-$m$ conference matrix $C$ in symmetric or skew form,
$C\otimes F$ is a weighing matrix of order $2m$ and weight $2m-2$.
If $C$ is skew, use a constant symmetric $R_i\in\mathcal D$;
if $C$ is symmetric, use a constant skew $R_i\in\mathcal D$.
The cross equations then vanish and produce a same-order Hadamard
matrix. These are familiar conference constructions, not new orders.
Exact checks find respectively eight fillings at order four,
six for the skew-conference tensor at order eight, and two for
the symmetric-conference tensor at order twelve.

The results separate two tasks: constructing a near-full weighing
matrix, and ensuring its small missing support has compatible
orthogonal-design structure. One zero per row has the classical
positive conference answer at admissible orders. Two zeros per
row permit a complete polynomial test, but the explicit counterexample
shows that existence of such a weighing matrix does not itself
force zero-preserving completion. Nothing here proves an all-order
Hadamard existence theorem or excludes repairs that also modify
some nonzero entries.
\section{Rank-two Gram defects: quantization and exact one-entry rounding}
\label{sec:spectral-rank-two}
Let $H\in\{-1,1\}^{n\times n}$, $4\mid n$, and
$E=HH^T-nI$. We study real rank, not binary rank. Small-rank Gram
perturbations are not assumed to preserve the sign constraint under
polar decomposition. The results below instead exploit the integer
square determinant and parity. Switching is classical background; see
Orrick, \emph{Switching operations for Hadamard matrices},
\url{https://arxiv.org/abs/math/0507515}. The elementary rectangular
flip calculation is included explicitly. Priority is unestablished.

\begin{theorem}[Complete energy quantization at rank two]
\label{thm:spectral-quantization}
If $E\ne0$ and $\operatorname{rank}_{\mathbb R}E=2$, there is an integer
$A\in\{1,\ldots,n/2\}$ with
\[
 \|E\|_F^2=8A(n-A),\qquad
 |\det H|=n^{(n-2)/2}(n-2A).
\]
The two nonzero eigenvalues of $E$ are
$\pm2\sqrt{A(n-A)}$. In particular the smallest possible nonzero
rank-two defect is $8(n-1)$. If $H_n$ exists, every listed value is
attained by flipping $A$ entries in one of its rows.
\end{theorem}
\begin{proof}
The zero diagonal gives trace zero, so the two nonzero eigenvalues are
$\lambda,-\lambda$, with $0<\lambda\le n$ by positivity of $HH^T$.
Also $\lambda^2=\sum_{i<j}E_{ij}^2\in4\mathbb Z$, because every
row inner product is even. Taking determinants gives
\[
 (\det H)^2=n^{n-2}(n^2-\lambda^2).
\]
Thus the nonnegative integer $n^2-\lambda^2$ is a rational square:
its square root is $|\det H|/n^{(n-2)/2}$. A rational number whose
square is an integer is itself an integer; call it $d$. Since
$d^2\equiv0\pmod4$, $d$ is even. Set $A=(n-d)/2$.
Then $1\le A\le n/2$ and $\lambda^2=n^2-d^2=4A(n-A)$.
For the attainment statement, keep $n-1$ Hadamard rows fixed and replace
the last by a version with $A$ signs flipped. The defect is a nonzero
star of rank two. Orthogonal expansion in the Hadamard row basis gives
$\lambda^2=n^2-(n-2A)^2=4A(n-A)$.
\end{proof}

\begin{theorem}[Energy-only one-entry repair certificate]
\label{thm:spectral-one-entry}
If any sign matrix of order $4\mid n$ satisfies
$\|HH^T-nI\|_F^2=8(n-1)$, it has a unique Hadamard neighbor obtained
by changing one entry. This statement needs no rank hypothesis.
Consequently a matrix with $\operatorname{rank}E\le2$ and
$\|E\|_F^2\le8(n-1)$ is either already Hadamard or admits this exact
one-entry repair.
\end{theorem}
\begin{proof}
Partition rows by parity of their numbers of negative entries, with
class sizes $a,b$. Same-parity inner products are divisible by four;
opposite-parity inner products are congruent to two modulo four.
If all rows have the same parity, $\|E\|_F^2$ is divisible by $32$,
whereas $8(n-1)$ is not. Hence $a,b>0$. The cross-parity pairs alone
contribute at least $8ab\ge8(n-1)$. Equality forces
$\{a,b\}=\{1,n-1\}$, every cross correlation to be $\pm2$, and all
other correlations zero. Thus $E$ is a star.

Delete its center row $v$ to obtain $(n-1)$ orthogonal sign rows $R$.
The matrix $P=nI-R^TR$ is positive semidefinite of rank one with diagonal
one. It is therefore $uu^T$ for a sign vector $u$, and $Ru=0$.
Orthogonal expansion gives
\[
 (v\cdot u)^2=n^2-\|Rv\|^2=n^2-4(n-1)=(n-2)^2.
\]
Choose the sign of $u$ to make $v\cdot u=n-2$; then $v,u$ differ in
exactly one entry. Replacing $v$ by $u$ repairs the matrix. This is
an integer algorithm: take $u_i=P_{i0}$, then choose its overall sign.
No other row can support a one-entry repair, because a nonzero star
edge avoiding that row would remain unchanged. Within the center row,
its orthogonal complement is the line $\mathbb Ru$, whose only sign
vectors are $\pm u$, and only one is at distance one. This proves
uniqueness. The last assertion follows from rank-two quantization;
a nonzero rank-one symmetric zero-diagonal matrix does not exist.
\end{proof}

\begin{lemma}[Rectangular flips explain, but do not classify, rank two]
\label{lem:spectral-rectangle}
Suppose a Hadamard matrix $H_0$ contains an $r\times c$ rank-one sign
rectangle, with sign vectors $u,v$ supported on its row and column sets.
Flip that rectangle to obtain $H=H_0-2uv^T$, and put $A=rc$.
Then $A\le n$. If $0<A<n$, the Gram defect has rank two and energy
$8A(n-A)$; if $A=n$ the flip preserves Hadamard orthogonality.
\end{lemma}
\begin{proof}
On the selected rows, $H_0v=cu$. The defect is
$uw^T+wu^T$, where $w=-2H_0v+2cu$ vanishes on those rows.
Thus $u,w$ have disjoint supports. Since $\|H_0v\|^2=nc$,
$\|w\|^2=4c(n-rc)$ and $\|u\|^2=r$. Positivity yields $rc\le n$,
and the two nonzero eigenvalues are
$\pm\|u\|\|w\|=\pm2\sqrt{A(n-A)}$ when $A<n$.
\end{proof}

\begin{proposition}[Rank two does not ensure a one-row or one-column repair]
\label{prop:spectral-nontrivial-repair}
There are explicit order-eight and order-twelve matrices with rank-two
Gram defect that require at least two rows and two columns to be altered
in any repair. Their minimum total number of entry changes is exactly
four. The order-twelve example is nonsingular.
\end{proposition}
\begin{proof}
In Walsh order eight, or the normalized Paley order twelve matrix, the
rectangle on rows $0,1$ and columns $0,2$ is all positive. Flip it.
The row defect graph is respectively $K_{2,2}$ or $K_{2,4}$, with all
nonzero correlations of magnitude four; the column defect graph has
the same respective form. Any unchanged nonorthogonal row or column
pair prevents repair, and both graphs have minimum vertex cover two.
To lower-bound total entry changes, let $f_i$ count changes in row $i$.
Every initially nonorthogonal pair needs
$2(f_i+f_j)\ge|E_{ij}|$. Summing over $K_{2,b}$ gives
$2b\sum_i f_i\ge 8b$, hence at least four changes. Reversing the
rectangle attains four. The order-twelve example has $A=4$, so its
nonzero determinant is $12^5(12-8)=4\cdot12^5$.
The checker verifies every stated graph and repair directly.
\end{proof}

More generally the same argument gives a global repair lower bound:
for arbitrary nonnegative edge weights $w_{ij}$ with positive maximum
weighted degree $\Delta$, every Hadamard repair changes at least
\[
 \frac{\sum_{i<j}w_{ij}|E_{ij}|}{2\Delta}
\]
entries. This is a concrete defect-graph budget certificate, not a
sufficiency theorem. Rank-two quantization and the one-entry threshold
provide a certified local rounding case; beyond it neither rank two nor
real polar orthogonalization alone establishes a flat repair.

\paragraph{The underlying defect graph.}
Every real symmetric rank-two matrix with zero diagonal has form
$uv^T+vu^T$ with disjoint nonempty supports for $u,v$. Indeed write
$E=\lambda(xx^T-yy^T)$ using orthonormal eigenvectors. The zero diagonal
gives $x_i^2=y_i^2$; the vectors
$u=\sqrt{\lambda/2}(x+y)$ and
$v=\sqrt{\lambda/2}(x-y)$ have disjoint supports and give the factorization.
Thus the nonzero-defect graph is complete bipartite, possibly with isolated
vertices. This structural fact does not say that every such sign matrix
is a single rectangular flip of a Hadamard matrix; no such classification
or general repair theorem is claimed.
\section{Flat real POVMs: exact commuting rigidity and a descent obstruction}
For an unbiased pairwise-independent law on $n-1$ sign bits write
$v_x=(1,x)\in\{\pm1\}^n$ and keep only positive-weight atoms.
Then $\sum_xw_xv_xv_x^T=I_n$. Set
\[
 u_x=v_x/\sqrt n,\quad P_x=u_xu_x^T,\quad
 E_x=nw_xP_x=w_xv_xv_x^T,\quad\lambda_x=nw_x.
\]
The $E_x$ form a real rank-one positive operator-valued measurement
(POVM): they are positive semidefinite and sum to identity.
Here ``flat'' means that the unit vectors have entries of magnitude
$1/\sqrt n$, not merely that the effect weights are equal.
Distinct atoms with the first coordinate normalized positive give
distinct nonparallel directions. The POVM inequality $E_x\le I_n$
implies $0<\lambda_x\le1$.

\begin{theorem}[Exact projective and commuting criteria]
For such a distinct-direction flat POVM the following are equivalent:
the law is a uniform Hadamard law; every $E_x$ is a projection;
the effects commute pairwise; and
$\sum_{x,y}\|[E_x,E_y]\|_F^2=0$.
Its total projectivity defect is exactly
\[
 \sum_x\operatorname{tr}(E_x-E_x^2)=n-n^2C(w).
\]
Writing $s_{xy}=v_x\cdot v_y$ and $A=\sum_xE_x^2$, its ordered
commutator energy is
\[
 \mathcal J=2\sum_{x,y}w_x^2w_y^2s_{xy}^2(n^2-s_{xy}^2)
 =2\left[\operatorname{tr}(A^2)-
                 \sum_{x,y}(\operatorname{tr}(E_xE_y))^2\right].
\]
\end{theorem}
\begin{proof}
The eigenvalue of $E_x$ on its image is $\lambda_x$, so
$\operatorname{tr}E_x=nw_x$ and $\operatorname{tr}E_x^2=n^2w_x^2$.
Two nonzero rank-one symmetric effects commute precisely when their
directions are orthogonal or parallel: their products are scalar
multiples of $u_xu_y^T$ and $u_yu_x^T$. Parallelism is excluded here.
Thus commuting forces mutually orthogonal directions. Completeness
then requires exactly $n$ of them, each with eigenvalue one, giving
the normalized rows of a Hadamard matrix. Conversely those rows give
commuting projections. A sum of squared norms vanishes exactly when
each commutator does. Finally expand
$[E_x,E_y]=w_xw_ys_{xy}(v_xv_y^T-v_yv_x^T)$ and use
$\|v_x\|^2=\|v_y\|^2=n$ to obtain the energy formula.
\end{proof}
Aggregating duplicate directions is essential to the commuting
criterion: splitting one projector into several parallel effects
would otherwise give a commuting but nonprojective measurement.

\begin{theorem}[Uniform tight frames are uniformly almost commuting]
If the law is uniform on $m$ atoms, put $\lambda=n/m$. Then
\[
 \max_x\|E_x-E_x^2\|_{\rm op}=\lambda(1-\lambda),\qquad
 \max_{x,y}\|[E_x,E_y]\|_{\rm op}^2
 \le\lambda^3(1-\lambda),\qquad
 \mathcal J/n\le2\lambda^2(1-\lambda).
\]
For every positive integer $r$, the normalized projectivity defect
of the self-tensor POVM is $1-(nC(w))^r$.
No self-tensor power becomes projective or commuting unless the
original POVM was already a Hadamard measurement.
\end{theorem}
\begin{proof}
Write $c_{xy}=u_x\cdot u_y$. Tightness gives
$\lambda\sum_yc_{xy}^2=1$, and hence
$\sum_{y\ne x}c_{xy}^2=(1-\lambda)/\lambda$.
The rank-two commutator has operator norm
$\lambda^2|c_{xy}|\sqrt{1-c_{xy}^2}$ and squared Frobenius norm
$2\lambda^4c_{xy}^2(1-c_{xy}^2)$. Bounding $1-c_{xy}^2\le1$
and summing proves both estimates. Tensor collision is $C(w)^r$,
so the defect identity gives the displayed normalized value.
If some base commutator is nonzero, fix any other nonzero effect $E_z$.
The tensor effects with first factors $E_x,E_y$ and the remaining
factors $E_z$ have commutator
$[E_x,E_y]\otimes(E_z^2)^{\otimes(r-1)}\ne0$. This excludes
commuting activation, and projectivity is excluded by the defect.
\end{proof}

\subsection{Even flat dilations with vanishing overhead do not ensure descent}
The finite-dimensional real dilation is explicit. The $m\times n$
matrix $W$ with rows $\sqrt{w_x}v_x^T$ satisfies $W^TW=I_n$, and
$E_x=W^T|x\rangle\langle x|W$. Complete its orthonormal columns to
a real orthogonal matrix. This is the elementary finite-rank
Naimark construction; the added columns have no automatic flatness.

In fact the character construction established in the earlier
stability branch gives a stronger counterexample. Use $h=2^s$,
known Hadamard order $k<h$, $g=h-k\ge s$, and slopes
$T\subset\mathbb F_h$ of size $k$ containing zero, whose omitted
set spans $\mathbb F_h$ and has size $g<h/2$.
For seeds $(a,b)$, the observed blocks are all nonconstant characters
of $a+tb$, $t\in T$, plus the mixed characters $\chi_e(tb)$ for
$t\ne0$. The resulting flat real POVM has
$m=h^2$, $n=kh$, and $\lambda=k/h$.
All augmented columns are distinct Walsh characters on the two
seeds. Thus $W$ is literally $n$ columns of a normalized real Walsh
matrix of order $m$: an \emph{exact flat real dilation already exists}.
Its dimension overhead is $m/n-1=h/k-1$.

The source orthogonality graph is exactly $h$ disjoint cliques of
size $h$, and completing these cliques using a known order-$k$
Hadamard matrix gives
\[
 \inf_{q\in\mathcal H_n}\|w-q\|_{\rm TV}=1-1/k,
\]
where $\mathcal H_n$ denotes mixtures of normalized real Hadamard
laws. Along the established known-order subsequences $k/h\to1$ and
$g\to\infty$, both effectwise defects, the normalized commutator
energy and the relative dilation overhead tend to zero, while this
distance tends to one. Therefore there is no dimension-independent
flat-real projectivization estimate tending to zero from these
defects. The counterexample even starts with an explicit flat
Naimark dilation, so flat dilation alone cannot supply the descent.

The distance here uses the common outcome labels given by actual
normalized sign directions. On the maximally mixed state $I_n/n$,
the measurement outcome probabilities are exactly $w_x$. Thus the
same lower bound applies to any measurement distance that dominates
the outcome TV distance for that input with these labels. No claim
is made for arbitrary relabeling, postprocessing, or unrestricted
complex projective measurements. This distinction matters: the
dimension-deficient Naimark and POVM-simulation literature, for
example Kotowski--Oszmaniec, allows different projective directions
and noise/postprocessing; see
\url{https://arxiv.org/abs/2501.09339}. Their usage of nearly flat
effects concerns weights, not the entrywise flatness required here.
The present identities and classical dilation are not claimed new;
the quantitative obstruction inherits the earlier sign-law theorem
and does not prove missing Hadamard existence.

\section{The remaining existence steps}
The quaternary route requires an actual four-phase construction at every
even order, beyond the Fourier-block families excluded here. This is a
sufficient stronger target, not an assumed equivalence with all real
Hadamard existence. The weighing route requires compatible global
block data: even two zeros per row do not force completion of a
prescribed matrix. The Gram-error result is a local repair theorem;
we have no procedure reaching its threshold in every admissible order.
The quantum formulation does not remove the requirement of entrywise
real flatness, and unrestricted measurement simulation is a different task.

Useful next tests are to construct quaternary families outside the
Fourier-block orbit, classify the two-zero triangle constraints by
signed-permutation holonomy, and find a globally convergent sign-matrix
operation that reaches the certified repair energy without a target
Hadamard matrix as an input. No such all-order mechanism is proved here.
No Lean code, git commits, publication or submission is part of this pass.
\end{document}
